\section{Related Work}
\label{sec:rel-work}

% We now discuss related work.


\minihead{Cybersecurity and AI}
Recent work in the intersection of cybersecurity and AI falls in three broad
categories: human uplift, societal implications of AI, and AI agents.

In this work, we focus on AI agents and cybersecurity. The closest works to ours
shows that ReAct-style AI agents can hack ``capture-the-flag'' toy websites and
vulnerabilities when given a description of the vulnerability \cite{fang2024llm,
fang2024llm2}. However, these agents fare poorly in the zero-day setting. In
particular, it is challenging for agents to backtrack after exploring a dead
end. We show in our work that teams of AI agents can autonomously exploit
zero-day vulnerabilities. Our findings are of broader relevance to the 
community, as governmental agencies \cite{us-aisi,uk-aisi}, industrial labs 
\cite{weidinger2024holistic,anthropic}, and other parties 
are interested in measuring cybersecurity capabilities of AI agents.

The human uplift setting focuses on using AI (typically LLMs) to aid humans in
cybersecurity tasks. For example, recent work has shown that LLMs can aid humans
in penetration testing and malware generation \cite{happe2023getting,
hilario2024generative}. This work is especially important in the setting of
``script kiddies'' who deploy malware without special expertise. Based on this, 
and the work on AI agents, researchers have also speculated on
societal implications of AI on cybersecurity \cite{lohn2022will,
handa2019machine}.


\minihead{AI agents}
AI agents have becoming increasing powerful and popular. Recent, highly capable
AI agents are largely based on LLMs \cite{yao2022react, weng2023agent} and can
now perform tasks as complex as solving real-world GitHub issues
\cite{yang2024sweagent}. There have been hundreds of papers on improving AI
agents, ranging from prompting techniques \cite{wei2022chain, yao2024tree},
planning techniques \cite{shinn2024reflexion, liu2023chain}, adding documents
and memory \cite{nuxoll2012enhancing}, domain-specific agents
\cite{he2024webvoyager}, and many more \cite{parisi2022talm}. The field 
of multi-agent systems is particularly related to our work \cite{liu2023dynamic,
chen2023autoagents, zhang2023building}. However, to the best of our knowledge,
our work is the first to introduce a real-world AI agent system based on
hierarchical planning and task-specific agents.


\minihead{Security of AI agents}
A related area of work is the security of AI agents themselves
\cite{greshake2023more, kang2023exploiting, zou2023universal, zhan2023removing,
qi2023fine, yang2023shadow}. Deployers of AI agents may want to limit the tasks
that the AI agent can do (e.g., restricting the ability to perform cybersecurity
attacks) and protect the agent against malicious attackers. Unfortunately,
recent work has shown that it is simple to bypass protections in LLMs, such as
by fine-tuning away protections \cite{zhan2023removing, yang2023shadow,
qi2023fine}. AI agents can also be attacked via indirect prompt injection
attacks \cite{greshake2023not, yi2023benchmarking, zhan2024injecagent}. This
line of work is orthogonal to ours.
